\section{Introduction}

Coastal flooding causes approximately \$1 trillion in global damages annually and
displaces hundreds of thousands of people~\citep{yang2023storm}. Reliable forecasting
of extreme water levels during storm events is therefore an urgent public safety
problem. The core challenge is \emph{data imbalance}: the most operationally
critical storm surges are also the rarest. Training standard neural forecasters on
historically recorded water levels allocates most model capacity to predicting the
frequent benign conditions rather than the infrequent but catastrophic extremes.

Data augmentation has emerged as a promising approach to rebalance training
distributions. For time-series forecasting, augmentation traditionally takes the
form of oversampling, bootstrapping, noise perturbation, or interpolation. These
methods are well-established for image data but have received less scrutiny for
coastal water-level prediction, where physical laws (tidal cycles, wave dynamics,
storm surge propagation) impose strong constraints on what constitutes a plausible
augmented sample.

Conditional diffusion models~\citep{ho2020denoising} offer a generative framework
that could produce physically realistic extreme events by learning the conditional
distribution of water-level residuals given historical gauge observations. The
hypothesis driving this work is:

\begin{quote}
\textbf{H1}: Tail-aware conditional diffusion augmentation, applied under a matched
forecaster training budget, improves prediction accuracy on \emph{independent held-out
real storm events} relative to real-only training, standard oversampling, perturbation
augmentation, and standard conditional diffusion.
\end{quote}

A secondary hypothesis (H2) investigates whether adding physical consistency or
multi-gauge spatial coherence further improves synthesis usefulness. Both are
\emph{unverified hypotheses} at this writing. This paper reports the frozen
research protocol, the engineering implementation, and the first controlled
empirical evidence — which is, at this stage, a \emph{negative result}: our
pilot DDPM produces samples of physically implausible magnitude. We argue that
this finding, its diagnosis, and the methodological infrastructure built to obtain
it are themselves a contribution to the study of diffusion models for extreme-event
time-series generation.

\subsection{Contributions}
\begin{enumerate}
    \item We present a fully documented, leakage-safe research protocol for
    coastal water-level forecasting with diffusion augmentation, including
    event-grouped chronological splits, full-support embargo, train-only thresholds,
    and a six-region data manifest frozen at license-admission boundaries.
    \item We implement and audit a conditional denoising diffusion probabilistic
    model (DDPM) for residual water-level generation, with 27,169 parameters, a
    cosine noise schedule, and 64 reverse diffusion steps.
    \item We report a \emph{reproduced negative result}: independently verified on
    both CUDA and CPU, the model's 24-hour generated residuals span $-16.91$ to
    $+7.94$ metres (SD 4.47m), approximately $22.95\times$ the training target
    standard deviation of 0.215m. We diagnose terminal-step amplification and
    imperfect epsilon prediction as contributing factors.
    \item We present engineering validation guards (nonfinite rejection,
    train-reference scale screen) that prevent absurdly bad samples from being
    misclassified as successful generation. These guards are validation
    infrastructure, not model repairs.
\end{enumerate}

\subsection{Paper structure}
Section~\ref{sec:related} surveys related work. Section~\ref{sec:protocol} details
our frozen research protocol and data admission. Section~\ref{sec:method} describes
the conditional DDPM architecture and training setup. Section~\ref{sec:exp}
describes our experimental plan and implementation.
Section~\ref{sec:results} reports our empirical findings, including the negative
result of the pilot. Section~\ref{sec:discussion} discusses implications, limitations,
and ongoing work. Section~\ref{sec:conclusion} concludes.
